% Separate post-CP19 proved source. Controlled reference data, not a count by itself.
\subsection{Controlled references for the actual one-curve coefficient gerbe}

\begin{lemma}[Proved: controlled one-curve gerbe anchors]
\label{D:controlled-one-curve-gerbe-anchors}
Use the setting of Proposition~\ref{D:one-curve-coefficient-root-gerbe}.
Fix an ordinary-$C$ depth $a$, and fix the characteristic $p$.
Let $n=p^f-1$, let $\alpha\in H^2(U,\mu_n)$, and let $m\mid n$
be its geometric Brauer period. Put
$\mathcal X=\sqrt[m]{(Y,C)}$ and let $C'\to C$ be the connected
degree-$m$ residue cover; $m=1$ is allowed.

The actual coefficient gerbe, its scalar-weight anchors at the
$f$ Frobenius stages, and their descent and comparison maps can
be chosen with the following properties.
\begin{enumerate}
\item All reference lines on $C'$ that enter a single root
layer of a termwise Hom complex, or a one-step Frobenius
comparison, have degree bounded in absolute value by $c m$.
Here $c$ depends only on $(Y,C),a,p$ and the fixed finite
diagram pattern. Root characters are reduced modulo $m$.
\item Through root depth $am$, the scalar-zero transition,
descent, and comparison cochains have rational representatives
whose residual pole divisors and plane-coordinate degrees are
bounded by a polynomial in $m$. This assertion concerns chosen
references, not arbitrary formal trivializations.
\item On adding the moving determinant lines of
Proposition~\ref{D:moving-line-polynomial-background}, the
reference data meet that proposition's polynomial background
hypotheses. For diagrams of bounded adjacent-phase range, the
polynomial is independent of $f$ and of $n$. The resulting
identities hold over arbitrary coefficient algebras on the
specified rational charts. No new moduli coordinates are needed
for the fixed reference cochains.
\end{enumerate}
In particular the normal-layer cohomology budget remains
$O(ms^2)$ per layer for spectral rank $s$, even though its
rational coefficient complexity may be a larger polynomial
in $m,s$. The lemma does not assert a representation count
without the effective bundle and unit-Isom presentations.
\end{lemma}

\begin{proof}
We distinguish the degree of a line from the pole degree of a
rational presentation. The former controls cohomology; the
latter controls the coefficients of the polynomial diagrams.
Constants below may increase with the fixed ordinary depth.

\emph{Separate the scalar order from the geometric period.}
Choose a line $H$ on $Y$ with $\deg(H|_C)=1$, using the
Picard lifting in Lemma~\ref{D:local-Picard-prime-to-p-generators}.
Choose $\gamma\in H^2(U,\mu_m)$ mapping to the Brauer image
$\beta$ of $\alpha$. In the Kummer diagram for $m\mid n$,
the map on the middle $\mathbf G_m$ is the identity and the
map on the right $\mathbf G_m$ is the $(n/m)$th power.
Thus pushout along $\mu_m\hookrightarrow\mu_n$ preserves
the Brauer image. The Picard calculation in
Proposition~\ref{D:one-curve-coefficient-root-gerbe} gives
\[
 \alpha=\iota_*\gamma+\delta_n(H^e|_U),\qquad
       0\leq e<\gcd(n,d)\leq d.                         \tag{D.CA.1}
\]
Here and below an empty choice of a nonzero character when
$m=1$ means the identity cover. Formula (D.CA.1) retains
the actual coefficient class. It does not replace it by its
Brauer image. Kummer exactness is the usual one for invertible
orders \citep[Lemma 59.28.1]{StacksControlledAnchorsD}.

Extend $\gamma$ to the $m$th root stack, as in the cited
coefficient-gerbe proposition. Represent the second summand
by the $n$th-root gerbe of $H^e$. If $T$ is its tautological
line, put $u_i$ equal to the residue of $p^i$ in $[0,n-1]$.
Then
\[
 (T^{u_i})^p=T^{u_{i+1}}\otimes H^{e c_i},\qquad
 c_i=\left\lfloor\frac{p u_i}{n}\right\rfloor,
                         0\leq c_i\leq p-1.             \tag{D.CA.2}
\]
There is no ordinary line of degree proportional to $n$ in
this comparison. Scalar-zero Hom cancels the tautological
line; the only ordinary one-step contribution is the bounded
power in (D.CA.2). Its pullback degree on $C'$ is $O(m)$.
All subsequent constructions keep this factor separate.

\emph{Controlled geometry of the residue cover.}
Fix a separating finite map $x:C\to\mathbb P^1$ once.
Its composite on $C'$ is separable of degree $O(m)$, and
$g(C')-1=m(g(C)-1)$. The plane-coordinate construction in
Proposition~\ref{D:moving-line-polynomial-background} therefore
gives a primitive equation of degree $O(m)$, discriminant
degree $O(m^2)$, and rational expressions of polynomial degree
for functions with a polynomially bounded pole divisor.
Deck translations preserve $x$ and satisfy the same bounds.

For completeness the cyclic cover itself requires no
uncontrolled defining functions. Its character is a pair
$(P,P^m\simeq\mathcal O_C)$ with $P\in\operatorname{Pic}^0(C)$.
Riemann--Roch represents $P$ by a difference of two divisors
of degrees bounded on the fixed curve $C$; the chosen
power trivialization then has divisor degree $O(m)$.
All changes to the preceding plane chart are of polynomial
degree. The finitely many fixed poles can be avoided by
the moving divisors, as in the background proposition.

Choose a fixed finite affine cover of $C$ on which its
normal line is trivial. A formal neighbourhood over one
such smooth affine curve is a product with a formal disc:
lift the retraction successively by smoothness, and choose
a generator of the normal ideal; the affine obstruction
groups in positive cohomological degree vanish.
Only the first $a$ ordinary normal coefficients of these
fixed charts are used. Their poles are fixed constants.
Finite etale invariance lifts $C'$ uniquely to these charts;
on a product chart the lift is the corresponding product.
Across charts its transition is the unique etale lift of
the residual transition. In a primitive presentation the
coefficient recursion divides by the invertible residual
derivative. Its inverse has polynomial pole degree by
the discriminant bound. Since there are only $a$ ordinary
normal steps, all these coefficients have polynomial pole
degree in $m$. Equivalently this is the explicit infinitesimal
form of topological invariance of the etale site
\citep[Theorems 59.45.1--2 and Proposition 59.45.4]{StacksControlledAnchorsD}.

Let $Y'\to Y$ be the etale lift, and let $\mathcal X'$ be
the base change of $\mathcal X$; write $D'$ for its root
divisor, so $mD'=C'$ and the degree of $\mathcal O(D')$
on the residual root gerbe is $-d$. If $t=z^m$ on a root
chart, a change of normal coordinate is obtained from its
ordinary change by a unique root of a unit congruent to one.
Through $am$ root layers this still uses only $a$ ordinary
layers. Its degree bound is made explicit below.

\emph{A normalized period-$m$ tautological line.}
The pullback of the geometric Brauer class to $\mathcal X'$
is zero: it is zero on the puncture, and restriction of the
Brauer group of this regular root stack is injective there,
as used in Proposition~\ref{D:one-curve-coefficient-root-gerbe}.
Thus the pulled-back $\mu_m$-gerbe is the root gerbe of
a line $M$ on $\mathcal X'$. The already proved local
Picard calculation, now on $Y'$, and the root-character
description of lines show that its class modulo $m$ has
a representative
\[
       M=(H')^b\otimes\mathcal O(tD'),
                   0\leq b,t<m,\qquad \deg(H'|_{C'})=1.
                                                               \tag{D.CA.3}
\]
Indeed write a line as an ordinary line times a root
character, reduce that character modulo $m$, and use
$\operatorname{Pic}(Y')/m\simeq\mathbb Z/m$ for the
remaining ordinary line. The formal Picard kernel and
the degree-zero residual Picard group are $m$-divisible.
The line $H'$ has a controlled formal representative:
choose one point in a product chart of $C'$ and take its
constant transverse Cartier section, gluing with the
trivial divisor away from that point. Formal GAGA supplies
an algebraic representative if desired. Formula (D.CA.3)
therefore has both residual degree $O(m)$ and controlled
normal transitions. Let $\Lambda$ denote its tautological
weight-one line. This is a choice of a controlled root
gerbe model, not a bound imposed on an arbitrary extension
of a previously selected $\Lambda$.

\emph{Retain and lift the actual descent data.}
Let $\tau$ be a generator of the cyclic cover. On the
residual root gerbe retain the descent data induced by
the original period-$m$ gerbe. In the model (D.CA.3) this
consists of a ratio line $Q$, its power isomorphism
\[
                 Q^m\simeq\tau^*M\otimes M^{-1},       \tag{D.CA.4}
\]
and the actual cyclic coherence, in particular the specified
trivialization of the norm product of the translates of $Q$.
These exist because the original gerbe descends. They are
not replaced by an arbitrarily chosen trivial norm.
The sign of $Q$ in an alternative convention changes none
of the following estimates.

The line $Q$ can have a nonzero root character. Write
$Q=Q_0\otimes\mathcal O(jD')$ with $0\leq j<m$.
Its total rational degree is zero by (D.CA.4), whence
\[
                        \deg Q_0=jd=O(m).               \tag{D.CA.5}
\]
Choose divisor representatives of $Q_0$, and of any
ordinary residual comparison lines, with positive and
negative degrees $O(m)$, using Riemann--Roch on $C'$.
Power isomorphisms and norm coherences are rational
isomorphisms between these specified divisor lines;
their pole degrees are polynomial in $m$ (an $m$th power
and a product of $m$ translates suffice). For powers of
$Q$, first reduce the root character modulo $m$. Their
ordinary line again has degree $O(m)$ because the total
degree is zero, and it can be replaced by a divisor model
of that degree. The comparison with the unreduced product
has polynomial pole degree, at worst adding a polynomial
number of these divisor degrees. This controls all root
characters without a second cover trivializing $D'$.

Lift these divisor lines by constant transverse sections
on the product charts. Such a lift $Q_{\mathrm{ref}}$ need
not yet satisfy (D.CA.4). Its defect is a line trivial on
the residual root gerbe. On the formal Picard kernel,
multiplication by $m$ is an isomorphism: each successive
square-zero kernel is a quotient of the first cohomology
of an additive conormal sheaf, multiplication by $m$ is
invertible there, and compatible roots pass to the inverse
limit. Equivalently use Kummer and invariance of prime-to-$p$
cohomology under nilpotent thickenings. Global residual
units are constants and lift, so no residual-unit boundary
introduces $m$-torsion in this kernel.

Here is the cochain version, including its degree control.
Trivialize the residual defect using the retained power
isomorphism, and extend those trivializations on the affine
charts. Its normal transition functions are $1+u$, with
$u$ in the normal ideal. There is a unique $v\equiv1$
with $v^m=1+u$. If
$v=1+\sum_{j\geq1}v_jz^j$, coefficient comparison at
order $j$ has the form
\[
       m v_j=u_j-P_j(v_1,\ldots,v_{j-1}),               \tag{D.CA.6}
\]
where every monomial of $P_j$ has total normal weight $j$.
The recursion divides only by the unit $m$. No factorial
or binomial denominator is used, even when $j\geq p$.
Thus each $v_j$ is a polynomial in the input coefficients
of weighted degree at most $j$. The same statement holds
for unit inversion. If an initial coefficient has a
polynomial pole bound, the bound at order $j\leq am$ is
still polynomial. Correct $Q_{\mathrm{ref}}$ by this root
of its defect, with the inverse sign when necessary.

Uniqueness preserves all commutative descent identities.
For example, the product of two chosen unit roots is the
chosen root of the product. The norm line has trivial
$m$th power and is residually trivial, hence is trivial
in the formal Picard kernel. Adjust its trivialization
by the unique unit root congruent to one so that its
$m$th power is the prescribed telescoping trivialization.
The remaining cyclic coherence defects have $m$th power
one and reduce to one, and hence equal one. The same
argument on triple overlaps gives exact cocycles.
We have therefore lifted the \emph{actual} residual
cyclic descent, not merely its isomorphism classes of
ratio lines. Nilpotent invariance of $\mu_m$-cohomology
identifies this controlled formal gerbe with the original
one at every retained order; the constructions are
compatible as the order increases. The root corrections
are commutative line/unit corrections, with no assertion
about arbitrary matrix corrections.

All coefficient bounds just used are rational pole bounds
on a fixed two-affine cover of $C'$. The representatives
for the necessary additive conormal classes may be chosen
in the finite-pole Cech complexes of
Proposition~\ref{D:moving-line-polynomial-background}.
Their residual line degrees and genus are $O(m)$, so the
initial pole range is polynomial in $m$. Formula (D.CA.6)
and the compatible pole homotopy increase that range by
at most a polynomial through the retained $am$ layers.
Trace-dual conversion in that proposition turns a
polynomial pole divisor into polynomial plane-coordinate
degree. Consequently evaluating these fixed functions
at moving divisor points has polynomial degree too.
All moving points can avoid the finitely many fixed pole
fibres; the ordered-divisor construction still represents
every required line. This verifies the coefficient bound
for a chosen reference model, including its gluing data.

\emph{Frobenius comparisons and the full scalar band.}
For the period-$m$ factor put $v_i=p^i\bmod m$, with
$0\leq v_i<m$. Use the tautological powers $\Lambda^{v_i}$,
with reduced root characters and the preceding controlled
divisor models. A one-step comparison has ordinary factor
$M^{\lfloor p v_i/m\rfloor}$, where the exponent is at
most $p-1$. Its ordinary and root-character normalized
residual line degrees are $O(m)$. Translated powers of
the actual ratio $Q$ have the controlled representatives
just constructed. The one-step Frobenius raises residual
coordinates to their fixed $p$th powers; it does not
iterate a growing-degree rational formula through all
$f$ phases. Every phase uses its own reduced-character
reference. Hence the same polynomial bound applies to
each bounded-range adjacent-phase block, independently
of $f$.

To recover the full coefficient class, use the contracted
product of the period-$m$ and $n$th-root models in
(D.CA.1). On their product the line
$\Lambda^{v_i}\otimes T^{u_i}$ descends: the contracted
$\mu_m$ acts trivially since $v_i\equiv u_i\pmod m$.
The only extra ordinary Frobenius factor is (D.CA.2).
No degree-$n$ ordinary splitting cover or line is
introduced. Even on local trivializations, a power
$w^{u_i}$ of a unit $w\equiv1$ has, at normal order $j$,
monomials of total normal weight at most $j$ independently
of $u_i$; its numerical binomial coefficients are simply
constants in $k$. When a unit root is required, use
(D.CA.6) with the invertible integer $n$. Thus an arbitrarily
large scalar order cannot enter the pole bound. Scalar-zero
termwise Hom cancels these tautological factors, leaving
exactly the ordinary controlled factors described above.

The residual spectral translations and their Frobenius
comparison maps are those of the retained descent data.
In particular the final comparison after $f$ steps is
retained as prescribed by
Proposition~\ref{D:residual-twisted-Frobenius-transport};
it is not declared to be the identity merely because
$p^f\equiv1\pmod n$.

Finally each normal layer of a termwise Hom has a residual
line of degree $O(m)$ and genus $O(m)$, tensored with the
bounded number of spectral term bundles. Its cohomology
cost is the $O(ms^2)$ cost of the fixed-determinant
construction. The larger polynomial pole range only
affects its rational coefficient matrices. Adding the
moving determinant lines changes the background degree
to a polynomial in $m,s$, exactly as in
Proposition~\ref{D:moving-line-polynomial-background}.
All formulas are identities of line bundles, unit
cocycles, and rational matrices before coefficient
extension. They therefore remain identities over every
coefficient algebra on the stated charts. This proves
all assertions without replacing geometric existence
of an anchor by an unproved quantitative bound.
\end{proof}
